\section{Discussion}\label{discussion}

\subsection{Practitioner implications}\label{discussion-practitioner-implications}

The main practical implication of the results is that a practitioner should read the benchmark output at the corpus level before considering any aggregate summary. The three dataset reports show that one organization can reasonably make different embedding-model choices for different document collections under the same evaluation roster and protocol. In the frozen snapshot, \texttt{M1} is the preferred candidate for \texttt{education\_eng} on operational grounds---it is not separated from \texttt{M2} on retrieval quality by the per-corpus paired intervals, but answers queries in roughly one third of the time; \texttt{M5} is the quality-first candidate for \texttt{education\_tr} and \texttt{law\_tr}, where its advantage over the runner-up is separated at the query level; and \texttt{M3} provides the lower-latency Turkish compromise when response time matters more than that measured increment in retrieval effectiveness. The per-corpus uncertainty layer makes this distinction visible by showing where the observed quality evidence separates the leading candidates and where operational evidence instead becomes the practical differentiator. These are examples of how to interpret the generated reports, not permanent prescriptions for future corpora.

In practical reuse, the practitioner first configures a benchmark for the corpus, candidate roster, and application requirements, then holds those choices constant while comparing the candidates and inspecting the dataset-level report. If the initial report motivates a different model roster or alternative retrieval, chunking, query, or evaluation settings, the pipeline can be rerun under that revised configuration as a separate sensitivity experiment. The present paper demonstrates the first decision-support cycle under one selected setup; it does not measure the effect of changing those settings.

This reading also clarifies the role of the aggregate report. The aggregate view shows that \texttt{M5} has the strongest cross-dataset quality summary in the present snapshot, but that summary does not override the dataset-level evidence recorded above. The intended workflow is therefore preliminary corpus-specific decision support: it narrows the candidate set under an explicit protocol and exposes quality--speed tradeoffs, while leaving final deployment choice to the practitioner's corpus, tolerance for synthetic evidence, and operational constraints.

\subsection{Language- and domain-aware corpus partitioning}\label{discussion-language-domain-partitioning}

Language and domain remain important because they help determine when a local collection should be partitioned and benchmarked separately. The English and Turkish education corpora produce different model-family rankings in this snapshot, whereas the Turkish education and Turkish law/regulation corpora preserve the same leading shortlist. The lesson is not that either language or domain is an isolated cause. Each comparison involves one corpus per condition, and corpus identity includes source selection, document structure, chunk distribution, generated queries, and qrels composition. The safer implication is that observed difference and observed stability should both be measured rather than assumed.

For practitioners, this supports a partition-first interpretation of the pipeline. If a local repository contains multiple languages, regulatory genres, institutional policy documents, or structurally different document types, the workflow can be run on each materially distinct subset. Separate reports then show whether the same candidate remains suitable or whether the shortlist changes. In the present paper, language-associated variation appears in the education comparison, while domain-associated stability appears in the Turkish comparison; both outcomes are useful because they prevent an unsupported transfer from one corpus report to another.

\subsection{Relation to prior benchmark workflows}\label{discussion-prior-workflows}

The contribution of this paper is an integrated local benchmark workflow rather than a new embedding algorithm, retrieval metric, or general-purpose leaderboard. Prior RAG and retrieval-evaluation work emphasizes the importance of retrieval quality, synthetic queries, LLM judging, benchmark governance, and reproducibility controls \citep{Lewis_RAG_2021,Zhu_LLMIR_Survey_2025,thakur2025support,abdallah2025rerankarena,miyaji2025evaluating,wang2025reproducibility,friel2024ragbench}. The pipeline connects those concerns in the single practitioner-facing sequence of Table~\ref{tbl:method_pipeline_stages}. As set out in Section~\ref{rw-tooling}, this places the workflow upstream of answer-quality evaluation frameworks and alongside, rather than in place of, the metric and dataset libraries on which it depends.

This positioning is intentionally modest. The paper does not claim that the five evaluated models exhaust the current embedding landscape or that the selected protocol is preferable to all alternatives. Instead, it demonstrates how a practitioner can turn a local document collection into reusable benchmark artifacts and compare an accessible candidate roster under one controlled setup. The value lies in making the comparison repeatable, inspectable, and reportable for each corpus, especially when manual trials would otherwise mix preprocessing choices, query sets, judgment policies, retrieval geometry, and reporting criteria.

\subsection{Limitations and transfer limits}\label{discussion-validity-transfer-limits}

Several limits constrain how far the reported rankings should be transferred. The benchmark uses synthetic queries and same-model LLM judging: Gemini 2.5 Flash-Lite both generates the queries and judges candidate relevance. Using the same model for both stages may introduce correlated model-specific biases between query construction and relevance assessment; the present study does not disentangle this dependency through an independent judge or human comparison. The paper includes illustrative examples and stored positive and negative labels, but it includes no human validation, no human--LLM agreement analysis, and no independent audit of query realism or label reliability. The pipeline deliberately does not require a pre-existing manually labeled query set or relevance judgments as input, because its intended operating scenario begins with an unlabeled local document collection. This design choice explains why human labels are not a prerequisite for running the workflow; it does not remove the separate need to validate the reliability of the generated benchmark. For higher-assurance use, a sampled human audit or judge-calibration stage can therefore be applied before model-selection decisions are finalized. The experiment generated 150 queries, but statistical scoring used 148 positive-bearing queries after zero-positive exclusion. The qrels are also roster-conditional because they are built from the union of candidates retrieved by the five evaluated models; relevant chunks outside the union of the pooled top-ranked candidates may remain unjudged \citep{zobel1998reliable, buckley2004retrieval}, and the qrels should not be treated as complete ground truth for arbitrary future systems.

The retrieval and protocol choices are similarly bounded. All rankings are conditional on default non-normalized Sentence-Transformers embeddings searched with ChromaDB squared-Euclidean (L2) distance; no normalized-cosine or model-recommended similarity rerun is reported. The evaluated roster contains five Sentence-Transformers candidates and excludes lexical retrieval, hybrid retrieval, reranking, domain-adapted models, newer multilingual models, and commercial alternatives. The selected chunk size, overlap, 50-query generation target, top-20 pool depth, and \texttt{k=10} cutoff are demonstration settings rather than tuned optima. Although the pipeline permits alternative configurations, functional configurability is not evidence that the reported rankings are robust to those alternatives; no sensitivity study is included.

Operational and reproducibility claims are also descriptive. The timing evidence combines corpus embedding with index insertion and combines query embedding with vector search; it comes from a single-run, single-machine, CPU-only snapshot and does not measure repeated-trial variability or runtime memory. The repository makes implementation and benchmark artifacts inspectable; external service behavior and model aliases may nevertheless change over time. These limits mean the pipeline output should be used as preliminary corpus-specific decision support, not as a deployment guarantee or as evidence that the observed ordering will transfer unchanged to another corpus, model roster, retrieval geometry, or hardware environment.
